\section{讨论与限制}

前述方法和实验应在本文的方法边界内理解。\SPMPC{} 的目标不是在规划器内部重建完整自由液面，也不是给出形式化无溢出保证；其核心作用是在可实时求解的移动底盘局部规划问题中保留液体动态记忆，使规划器能够在路径推进、控制平滑性和预测晃液响应之间进行权衡。

\subsection{低阶晃液模型的适用范围}

\SPMPC{} 使用低阶晃液模态状态，是为了在移动底盘局部规划的时间尺度内获得可用于优化的液体动态记忆。该选择适合表达主导晃液模态、历史激励和当前液体相位对未来响应的影响，也使晃液状态能够作为增强状态进入滚动时域 OCP。

这一建模选择也带来边界。低阶模型不能覆盖所有强非线性自由液面现象，例如破波、局部冲击、强三维效应、复杂容器几何、液位大幅变化以及外界扰动造成的非模型化液体运动。因此，本文中的晃液状态应理解为控制友好的预测状态，而不是高保真液体仿真的替代物。若面向更复杂容器、较高液位或剧烈运动任务，需要重新辨识模型参数，或引入更丰富的液体状态表示。

\subsection{模型代理量与真实液面评价}

本文必须区分两类量：规划器内部传播的模型预测晃液状态，以及实验中观测到的真实液面响应。前者包括低阶模态状态、预测晃液代价以及实现中可能发布的 \texttt{/spmpc/slosh\_height} 等诊断量；它们用于解释优化器为何选择某段运动，也用于检查模型预测是否符合预期。后者则应来自独立观测，例如 RGB 液面指标、液面轮廓变化、外部液位传感器或其它实验测量。

因此，内部晃液代理量不能直接等同于真实液面高度。若实验只显示模型预测量降低，只能说明规划器在自身模型下减小了预测晃液响应；只有当外部液面观测也出现一致趋势时，才能支持真实液体晃动降低的结论。这个区分对于实物实验尤其重要，因为相机视角、液面反光、容器边缘、传感延迟和标定误差都会影响真实液面指标。

\subsection{防溢出保证与安全边界}

本文不提供形式化的严格防溢出保证。即使在 OCP 中加入模型预测晃液边界或软约束，该约束限制的仍然是低阶模型状态，而不是真实自由液面的所有可能运动。在没有真实液面闭环反馈、保守集合构造和形式化安全证明的情况下，不能把模型预测边界表述为严格的无溢出保证。

因此，本文更合适的表述是降低晃动或降低溢出风险，而不是保证不溢出。若未来希望形成严格安全结论，需要进一步研究液面状态估计、模型不确定性界、约束收紧、鲁棒或机会约束形式，以及与 CBF 等安全机制的结合。

\subsection{局部规划层级与导航功能边界}

\SPMPC{} 的定位是标准轮式移动底盘上的在线局部规划器。本文假设已有安全参考路径，并在该路径附近优化路径进度、底盘命令、控制平滑性和预测液体响应。因此，本文不把完整全局路径搜索、动态障碍物交互、同伦路径选择、走廊生成、地图语义速度先验或复杂地形适应作为主贡献。

这一边界有助于保护论文主线。普通局部规划器、移动底盘防晃近邻工作和跨平台防晃控制文献各自解决了相关但不同层级的问题。本文的重点不是替代完整导航栈，也不是宣称在所有局部规划指标上优于所有方法，而是验证：当任务包含开口液体时，把液体动态记忆放入在线 \MPCC{} 局部规划层是否能产生仅靠平滑控制无法获得的作用。

\subsection{实物闭环与工程诊断边界}

在实物系统中，命令限幅、终端减速、传感延迟、底盘执行响应和状态估计噪声都会影响最终液面运动。因此，实物结果应同时记录求解器输出、最终底盘命令、外部液面观测和必要的执行链路诊断，以区分方法效果与工程实现影响。这些诊断量用于保证实验可解释性，但不构成 \SPMPC{} 的核心方法贡献；若未来将 IMU 或真实液面反馈纳入闭环，则需要重新定义状态估计问题和实验对照。

\subsection{未来工作}

未来工作可以从四个方向展开：引入真实液面反馈或容器加速度测量，形成闭环液体状态估计和参数自适应；将模型预测晃液边界扩展为更严格的防溢出约束，并分析模型误差和约束保守性；将晃液感知局部规划扩展到动态障碍、走廊约束、多路径选择或更复杂的导航场景；在不同液位、不同容器形状、多容器运输和真实实验室任务中评估方法鲁棒性。
